\scanpage{57}
\mspage{56}

On trouve

\uline{Scholie} des ateliers \inserted{(quasi)} ensemblistes à figures de type fini : se donner un tel atelier, revient au $=$ que de se donner un \deleted{\unsure{ens}} couple $(L,M)$, $L$ ensemble, $M$ ens. de figures élémentaires dans $L$, telles que toute sous figure $F$ d'un $X\in M$ soit de type fini (OK si les $X\in M$ sont des figures finies) et satisfaisant de plus At ens $\unsure{\sharp|\mathrm{gf}|{}^{**}}$ 1 -- At ens $\sharp$ 3.

On \unsure{réconstitut|réconstitue} $\unsure{\mathcal F|\mathscr F}$ comme sous-atelier de $\operatorname{Fig}(L)$, formé des figures \inserted{ens.} \uline{de type fini} telles que \unsure{les|ses} \unsure{multistrates} soient dans $M$.

Dans le cas d'un atelier ens. pas \unsure{néc.} \inserted{à fig.} de type fini, on \deleted{\unsure{préciserait|paraîtrait|présumerait} néanmoins que}
\deleted{pour une figure ensembliste $F$ dans $L$, dont les \deleted{strates} multistrates sont dans $M$, l'appartenance à $\mathcal F$ ne dépend que de son \uline{ombre}.}
